As we have indicated above, the key to obtaining an existence theory is proving an energy estimate analogous to Gårding\textquotesingle s inequality.

\begin{Theorem}
\label{thm:neumann-energy-estimate}
Let $L(\bm{x},D)$ be a second-order linear partial differential operator in divergence form of the form described in \eqref{eq:second-order-garding-divergence-operator} such that for some $\theta>0$ the uniform ellipticity condition \eqref{eq:second-order-ellipticity-negative-coefficient-matrix} holds. Also suppose that $a_{ij},b_k\in L^\infty(\Omega)$ for $i,j=1,\ldots,n$, $k=0,\ldots,n$. Then there exist constants $\overline c$ and $\lambda_N\geq0$ such that
\begin{equation}\label{eq:neumann-energy-coercivity-inequality}
B[u,u]+\lambda_N\|u\|_{L^2(\Omega)}^2\geq\overline c\|u\|_{H^1(\Omega)}^2\quad\text{for all }u\in H^1(\Omega),
\end{equation}
where $B$ is as defined in \eqref{eq:second-order-garding-bilinear-form}.
\end{Theorem}

\begin{Proof}
The statement of this theorem is identical to that of Theorem~\ref{thm:second-order-garding-inequality} except that now we are trying to prove the theorem over the space $H^1(\Omega)$ rather than just $H_0^1(\Omega)$. Thus, the only difference in the proof of this result is that we no longer have Poincaré\textquotesingle s inequality. However, as before we can get
\begin{equation}\label{eq:neumann-energy-gradient-lower-bound}
B[u,u]\geq\dfrac{\theta}{2}\int_\Omega|\nabla u|^2\,\mathrm{d}\bm{x}-C\|u\|_2^2.
\end{equation}
And instead of using Poincaré\textquotesingle s inequality at this point we simply write
\begin{equation}\label{eq:neumann-energy-full-h1-lower-bound}
B[u,u]\geq\dfrac{\theta}{2}\|u\|_{1,2}^2-\left(C+\dfrac{\theta}{2}\right)\|u\|_2^2,
\end{equation}
which completes the proof.
\end{Proof}

It is worth noting that the only real reason for using Poincaré in the proof of Theorem~\ref{thm:second-order-garding-inequality} was to get a sharper estimate on the constant $\lambda_G$. However, since we haven\textquotesingle t been trying to specify optimal constants anyway, this effort was sort of wasted.

With our energy estimate in place to take care of the coercivity condition in the Lax-Milgram lemma, the existence of a weak solution of the Neumann problems follows with only minor modifications of the proof of Theorem~\ref{thm:weak-dirichlet-positive-shift-existence}.

\begin{Theorem}
\label{thm:neumann-positive-shift-weak-existence}
Let $L(\bm{x},D)$ be a second-order linear partial differential operator in divergence form satisfying the hypotheses of Theorem~\ref{thm:neumann-energy-estimate}. Then there exists $\lambda_N\geq0$ such that for any $\widetilde\lambda\geq\lambda_N$ and for any $f\in L^2(\Omega)$, the Neumann problem for the operator
\begin{equation}\label{eq:neumann-positive-shifted-operator}
\widetilde L(\bm{x},D):=L(\bm{x},D)+\widetilde\lambda
\end{equation}
has a unique weak solution $u\in H^1(\Omega)$. Furthermore, this solution satisfies
\begin{equation}\label{eq:neumann-weak-solution-l2-data-bound}
\|u\|_{1,2}\leq C\|f\|_2.
\end{equation}

\end{Theorem}
